Set $x=\cot(\dfrac{A}{2})$, $y=\cot(\dfrac{B}{2})$, $z=\cot(\dfrac{C}{2})$. Since the half-angles sum to $\dfrac{\pi}{2}$, their tangents satisfy
\[
\dfrac{1}{xy}+\dfrac{1}{yz}+\dfrac{1}{zx}=1,
\]
or $x+y+z=xyz$. The arithmetic progression gives $x+z=2y$. Hence $3y=xyz$, and division by the positive number $y$ gives $xz=3$.
\begin{asy}
texpreamble("\usepackage{amsmath,amssymb}");
size(240);
defaultpen(fontsize(9pt));
pair A=(0,0),B=(4,0),C=(2,2sqrt(3));draw(A--B--C--cycle);label("$A$",A,SW);label("$B$",B,SE);label("$C$",C,N);label("Equal-angle example",(2,-.7));
\end{asy}
